% SOURCE_AUTHORITY: sga5_fr_workpass.tex, lines 13086--13109 (LF-normalized unit).
% SOURCE_SHA256: 791F4EFFC5E02832D5D77ED03518C8156D6F07E4C8238B03545DB93D883FBB28
\begin{statement}{Corolário 7.12}
Sejam $C$ uma curva algébrica conexa, lisa e própria sobre $k$, $V$ um aberto não vazio de $C$, $\Delta$ uma subcategoria de $D^+(A)$ como em 7.1, e $F$ um complexo de feixes de $A$-módulos sobre $V$ tal que sua restrição a um subaberto não vazio conveniente de $V$ seja trivializada por um recobrimento de Galois de $V$. Suponha-se ainda que $F_{\bar x}\in\Ob(\Delta)$ qualquer que seja $x\in V$. Nestas condições, tem-se:
\[
(*)\qquad \RG_{!V}(F),\ \RG_V(F)\in\Ob(\Delta),
\]
\[
(**)\qquad
\clK{\Delta}(\RG_{!V}(F))=\clK{\Delta}(\RG_V(F)).
\]
Ponhamos $\EP_\Delta(V,F)=\clK{\Delta}(\RG_V(F))$ e
\[
\EP(V)=\EP_{D(F_\ell)_{\coh}}(V,F_\ell)=\EP(C)-\operatorname{card}(Z),
\]
onde $Z=C-V$ e onde $F_\ell$ é o feixe sobre $C$ definido em dimensão zero pelo feixe constante associado a $F_\ell=\mathbb Z/\ell\mathbb Z$. Tem-se, então, a fórmula:
\[
\EP_\Delta(V,F)
=
\EP(V)\clK{\Delta}(F_{\bar\eta})
-\sum_{x\in V(k)}\varepsilon_x^\Delta(F)
-\sum_{x\in Z}\alpha_x^\Delta(F),
\tag{7.13}
\]
onde $\varepsilon_x^\Delta$ e $\alpha_x^\Delta$ estão definidos por (6.2.1), (6.2.2).
\end{statement}
